
\documentclass{article}
\usepackage{colortbl}
\usepackage{makecell}
\usepackage{multirow}
\usepackage{supertabular}

\begin{document}

\newcounter{utterance}

\centering \large Interaction Transcript for game `clean\_up', experiment `2\_hard\_3obj\_fr', episode 3 with glm{-}5.2.
\vspace{24pt}

{ \footnotesize  \setcounter{utterance}{1}
\setlength{\tabcolsep}{0pt}
\begin{supertabular}{c@{$\;$}|p{.15\linewidth}@{}p{.15\linewidth}p{.15\linewidth}p{.15\linewidth}p{.15\linewidth}p{.15\linewidth}}
    \# & $\;$A & \multicolumn{4}{c}{Game Master} & $\;\:$B\\
    \hline

    \theutterance \stepcounter{utterance}  
    & & \multicolumn{4}{p{0.6\linewidth}}{
        \cellcolor[rgb]{0.9,0.9,0.9}{
            \makecell[{{p{\linewidth}}}]{
                \texttt{\tiny{[P1$\langle$GM]}}
                \texttt{Je suis votre maître de jeu, et vous jouez à un jeu collaboratif avec la grille suivante comme plateau de jeu :} \\
\\ 
\texttt{```} \\
\texttt{ 1234567} \\
\texttt{╔═╤═╤═══╗} \\
\texttt{║L│◌│◌◌◌║ 1} \\
\texttt{║C├─┤◌◌◌║ 2} \\
\texttt{╟─┤◌├─┬─╢ 3} \\
\texttt{║◌│◌│◌│◌║ 4} \\
\texttt{╟─┼─┘◌│◌║ 5} \\
\texttt{║◌│◌◌┌┴─╢ 6} \\
\texttt{║◌│◌◌│◌P║ 7} \\
\texttt{╚═╧══╧══╝} \\
\\ 
\texttt{```} \\
\\ 
\texttt{* Le bord supérieur affiche des coordonnées x croissantes vers la droite, et le bord droit des coordonnées y croissantes vers le bas.} \\
\texttt{* Les objets suivants sont placés aléatoirement sur votre grille : 'C', 'L', 'P'.} \\
\\ 
\\ 
\texttt{L'autre joueur voit une variante du plateau de jeu, où les objets sont placés à des emplacements aléatoires différents. Vous ne pouvez pas voir le plateau de l'autre joueur et il ne peut pas voir le vôtre.} \\
\\ 
\\ 
\texttt{**Objectif :**} \\
\\ 
\texttt{Les deux joueurs doivent déplacer les objets sur leur arrière{-}plan respectif de manière à ce que des objets identiques se retrouvent aux mêmes coordonnées. Vous devez communiquer avec l'autre joueur pour convenir d'une configuration d'objectif commune.} \\
\\ 
\\ 
\texttt{**Règles :**} \\
\\ 
\texttt{* À chaque tour, vous pouvez envoyer exactement l'une des deux commandes suivantes :} \\
\texttt{ 1. `DIRE: <MESSAGE>` : pour envoyer un message (tout ce qui se trouve jusqu'au prochain saut de ligne) à l'autre joueur. Je le transmettrai à votre partenaire.} \\
\texttt{2. `DÉPLACER: <OBJET>, (<X>, <Y>)` : pour déplacer un objet vers une nouvelle position, où `<X>` est la colonne et `<Y>` est la ligne. Je vous informerai si votre déplacement a réussi ou non.} \\
\texttt{* Si vous ne respectez pas le format, ou si vous envoyez plusieurs commandes à la fois, je dois vous pénaliser.} \\
\texttt{* Si les deux joueurs accumulent plus de 8 pénalités, vous perdez tous les deux la partie.} \\
\texttt{* Il est essentiel que vous communiquiez avec l'autre joueur concernant votre état d'objectif ! La *seule* façon de transmettre votre stratégie à l'autre joueur est d'utiliser la commande `DIRE: <MESSAGE>` !} \\
\\ 
\\ 
\texttt{**Déplacement d'objets**} \\
\\ 
\texttt{* Vous ne pouvez déplacer des objets que vers des cellules situées à l'intérieur des limites de la grille. La cellule cible doit être vide, c'est{-}à{-}dire qu'elle ne doit contenir que le symbole '◌'.} \\
\texttt{* Si vous essayez de déplacer un objet vers un endroit qui n'est pas vide, ou si vous essayez de le déplacer en dehors de la grille, je dois vous pénaliser. Vous avez droit à un nouvel essai.} \\
\texttt{* Avant d'effectuer un déplacement, vérifiez bien que l'emplacement cible est vide et qu'il ne contient aucune lettre, aucun cadre ou aucune ligne !} \\
\\ 
\\ 
\texttt{**Fin de partie**} \\
\\ 
\texttt{Si vous pensez avoir atteint l'objectif d'aligner tous les objets, vous pouvez demander à l'autre joueur de terminer la partie en envoyant `DIRE: terminé ?`. Si l'autre joueur vous demande de terminer la partie et que vous répondez `DIRE: terminé !`, la partie se termine.} \\
\\ 
\texttt{Les deux joueurs gagnent si la partie se termine en 12 tours, un tour étant défini comme l'envoi d'une commande valide par chacun des deux joueurs.} \\
\\ 
\texttt{**Score :**} \\
\\ 
\texttt{Plus les objets identiques sont proches les uns des autres sur les deux plateaux de jeu, plus vous obtenez de points. Les pénalités réduisent vos points. Pouvez{-}vous battre le record ?} \\
\\ 
\\ 
\texttt{Envoyez un message à l'autre joueur pour commencer le jeu !} \\
            }
        }
    }
    & \\ \\

    \theutterance \stepcounter{utterance}  
    & \multicolumn{4}{p{0.6\linewidth}}{
        \cellcolor[rgb]{0.9,0.9,0.9}{
            \makecell[{{p{\linewidth}}}]{
                \texttt{\tiny{[P1$\rangle$GM]}}
                \texttt{DIRE: Bonjour ! Je propose qu'on aligne nos objets à (4,4), qui semble être une case vide centrale sur mon plateau. Mes objets sont actuellement en : L(1,1), C(1,2), P(7,7). Dites{-}moi vos positions de départ et si (4,4) vous convient comme cible commune.} \\
            }
        }
    }
    & & \\ \\

    \theutterance \stepcounter{utterance}  
    & & & \multicolumn{2}{p{0.3\linewidth}}{
        \cellcolor[rgb]{0.9,0.9,0.9}{
            \makecell[{{p{\linewidth}}}]{
                \texttt{\tiny{[GM$|$GM]}}
                \texttt{DIRE: Bonjour ! Je propose qu'on aligne nos objets à (4,4), qui semble être une case vide centrale sur mon plateau. Mes objets sont actuellement en : L(1,1), C(1,2), P(7,7). Dites{-}moi vos positions de départ et si (4,4) vous convient comme cible commune.} \\
            }
        }
    }
    & & \\ \\

    \theutterance \stepcounter{utterance}  
    & & \multicolumn{4}{p{0.6\linewidth}}{
        \cellcolor[rgb]{0.9,0.9,0.9}{
            \makecell[{{p{\linewidth}}}]{
                \texttt{\tiny{[GM$\rangle$P2]}}
                \texttt{Je suis votre maître de jeu, et vous jouez à un jeu collaboratif avec la grille suivante comme plateau de jeu :} \\
\\ 
\texttt{```} \\
\texttt{ 1234567} \\
\texttt{╔═╤═╤═══╗} \\
\texttt{║◌│◌│◌◌◌║ 1} \\
\texttt{║◌├─┤◌L◌║ 2} \\
\texttt{╟─┤P├─┬─╢ 3} \\
\texttt{║◌│◌│◌│◌║ 4} \\
\texttt{╟─┼─┘C│◌║ 5} \\
\texttt{║◌│◌◌┌┴─╢ 6} \\
\texttt{║◌│◌◌│◌◌║ 7} \\
\texttt{╚═╧══╧══╝} \\
\\ 
\texttt{```} \\
\\ 
\texttt{* Le bord supérieur affiche des coordonnées x croissantes vers la droite, et le bord droit des coordonnées y croissantes vers le bas.} \\
\texttt{* Les objets suivants sont placés aléatoirement sur votre grille : 'C', 'L', 'P'.} \\
\\ 
\\ 
\texttt{L'autre joueur voit une variante du plateau de jeu, où les objets sont placés à des emplacements aléatoires différents. Vous ne pouvez pas voir le plateau de l'autre joueur et il ne peut pas voir le vôtre.} \\
\\ 
\\ 
\texttt{**Objectif :**} \\
\\ 
\texttt{Les deux joueurs doivent déplacer les objets sur leur arrière{-}plan respectif de manière à ce que des objets identiques se retrouvent aux mêmes coordonnées. Vous devez communiquer avec l'autre joueur pour convenir d'une configuration d'objectif commune.} \\
\\ 
\\ 
\texttt{**Règles :**} \\
\\ 
\texttt{* À chaque tour, vous pouvez envoyer exactement l'une des deux commandes suivantes :} \\
\texttt{ 1. `DIRE: <MESSAGE>` : pour envoyer un message (tout ce qui se trouve jusqu'au prochain saut de ligne) à l'autre joueur. Je le transmettrai à votre partenaire.} \\
\texttt{2. `DÉPLACER: <OBJET>, (<X>, <Y>)` : pour déplacer un objet vers une nouvelle position, où `<X>` est la colonne et `<Y>` est la ligne. Je vous informerai si votre déplacement a réussi ou non.} \\
\texttt{* Si vous ne respectez pas le format, ou si vous envoyez plusieurs commandes à la fois, je dois vous pénaliser.} \\
\texttt{* Si les deux joueurs accumulent plus de 8 pénalités, vous perdez tous les deux la partie.} \\
\texttt{* Il est essentiel que vous communiquiez avec l'autre joueur concernant votre état d'objectif ! La *seule* façon de transmettre votre stratégie à l'autre joueur est d'utiliser la commande `DIRE: <MESSAGE>` !} \\
\\ 
\\ 
\texttt{**Déplacement d'objets**} \\
\\ 
\texttt{* Vous ne pouvez déplacer des objets que vers des cellules situées à l'intérieur des limites de la grille. La cellule cible doit être vide, c'est{-}à{-}dire qu'elle ne doit contenir que le symbole '◌'.} \\
\texttt{* Si vous essayez de déplacer un objet vers un endroit qui n'est pas vide, ou si vous essayez de le déplacer en dehors de la grille, je dois vous pénaliser. Vous avez droit à un nouvel essai.} \\
\texttt{* Avant d'effectuer un déplacement, vérifiez bien que l'emplacement cible est vide et qu'il ne contient aucune lettre, aucun cadre ou aucune ligne !} \\
\\ 
\\ 
\texttt{**Fin de partie**} \\
\\ 
\texttt{Si vous pensez avoir atteint l'objectif d'aligner tous les objets, vous pouvez demander à l'autre joueur de terminer la partie en envoyant `DIRE: terminé ?`. Si l'autre joueur vous demande de terminer la partie et que vous répondez `DIRE: terminé !`, la partie se termine.} \\
\\ 
\texttt{Les deux joueurs gagnent si la partie se termine en 12 tours, un tour étant défini comme l'envoi d'une commande valide par chacun des deux joueurs.} \\
\\ 
\texttt{**Score :**} \\
\\ 
\texttt{Plus les objets identiques sont proches les uns des autres sur les deux plateaux de jeu, plus vous obtenez de points. Les pénalités réduisent vos points. Pouvez{-}vous battre le record ?} \\
\\ 
\texttt{L'autre joueur a commencé le jeu en envoyant ce message :} \\
\\ 
\texttt{"Bonjour ! Je propose qu'on aligne nos objets à (4,4), qui semble être une case vide centrale sur mon plateau. Mes objets sont actuellement en : L(1,1), C(1,2), P(7,7). Dites{-}moi vos positions de départ et si (4,4) vous convient comme cible commune."} \\
\\ 
\texttt{Quelle est votre première commande ?} \\
            }
        }
    }
    & \\ \\

    \theutterance \stepcounter{utterance}  
    & & & \multicolumn{4}{p{0.6\linewidth}}{
        \cellcolor[rgb]{0.9,0.9,0.9}{
            \makecell[{{p{\linewidth}}}]{
                \texttt{\tiny{[GM$\langle$P2]}}
                \texttt{DIRE: Bonjour ! Sur mon plateau, (4,4) est libre et central aussi, ça me convient parfaitement comme cible commune. Pour information, mes positions de départ sont : C(4,5), L(4,2), P(2,3). Je commence par déplacer P vers (4,4).} \\
            }
        }
    }
     \\ \\

    \theutterance \stepcounter{utterance}  
    & & & \multicolumn{2}{p{0.3\linewidth}}{
        \cellcolor[rgb]{0.9,0.9,0.9}{
            \makecell[{{p{\linewidth}}}]{
                \texttt{\tiny{[GM$|$GM]}}
                \texttt{DIRE: Bonjour ! Sur mon plateau, (4,4) est libre et central aussi, ça me convient parfaitement comme cible commune. Pour information, mes positions de départ sont : C(4,5), L(4,2), P(2,3). Je commence par déplacer P vers (4,4).} \\
            }
        }
    }
    & & \\ \\

    \theutterance \stepcounter{utterance}  
    & & \multicolumn{4}{p{0.6\linewidth}}{
        \cellcolor[rgb]{0.9,0.9,0.9}{
            \makecell[{{p{\linewidth}}}]{
                \texttt{\tiny{[P1$\langle$GM]}}
                \texttt{Votre message a été transmis à l'autre joueur.} \\
\\ 
\\ 
\texttt{Vous jouez actuellement le tour 1 sur un maximum de 12.} \\
\texttt{Vous avez collectivement accumulé 0 pénalités sur 8.} \\
\texttt{L'autre joueur a envoyé ce message :} \\
\\ 
\\ 
\texttt{"Bonjour ! Sur mon plateau, (4,4) est libre et central aussi, ça me convient parfaitement comme cible commune. Pour information, mes positions de départ sont : C(4,5), L(4,2), P(2,3). Je commence par déplacer P vers (4,4)."} \\
\\ 
\\ 
\texttt{Quelle est votre prochaine commande ?} \\
            }
        }
    }
    & \\ \\

    \theutterance \stepcounter{utterance}  
    & \multicolumn{4}{p{0.6\linewidth}}{
        \cellcolor[rgb]{0.9,0.9,0.9}{
            \makecell[{{p{\linewidth}}}]{
                \texttt{\tiny{[P1$\rangle$GM]}}
                \texttt{DIRE: D'accord pour (4,4) comme cible commune. Comme vous occupez bientôt (4,4), je vais déplacer mon L en (4,4) ce tour{-}ci aussi. Ensuite, je vous laisse déplacer C ou L ailleurs, puis j'ajusterai mon C et P sur la nouvelle cible. Faites{-}moi savoir votre prochaine cible.} \\
            }
        }
    }
    & & \\ \\

    \theutterance \stepcounter{utterance}  
    & & & \multicolumn{2}{p{0.3\linewidth}}{
        \cellcolor[rgb]{0.9,0.9,0.9}{
            \makecell[{{p{\linewidth}}}]{
                \texttt{\tiny{[GM$|$GM]}}
                \texttt{DIRE: D'accord pour (4,4) comme cible commune. Comme vous occupez bientôt (4,4), je vais déplacer mon L en (4,4) ce tour{-}ci aussi. Ensuite, je vous laisse déplacer C ou L ailleurs, puis j'ajusterai mon C et P sur la nouvelle cible. Faites{-}moi savoir votre prochaine cible.} \\
            }
        }
    }
    & & \\ \\

    \theutterance \stepcounter{utterance}  
    & & \multicolumn{4}{p{0.6\linewidth}}{
        \cellcolor[rgb]{0.9,0.9,0.9}{
            \makecell[{{p{\linewidth}}}]{
                \texttt{\tiny{[GM$\rangle$P2]}}
                \texttt{Votre message a été transmis à l'autre joueur.} \\
\\ 
\\ 
\texttt{Vous jouez actuellement le tour 2 sur un maximum de 12.} \\
\texttt{Vous avez collectivement accumulé 0 pénalités sur 8.} \\
\texttt{L'autre joueur a envoyé ce message :} \\
\\ 
\\ 
\texttt{"D'accord pour (4,4) comme cible commune. Comme vous occupez bientôt (4,4), je vais déplacer mon L en (4,4) ce tour{-}ci aussi. Ensuite, je vous laisse déplacer C ou L ailleurs, puis j'ajusterai mon C et P sur la nouvelle cible. Faites{-}moi savoir votre prochaine cible."} \\
\\ 
\\ 
\texttt{Quelle est votre prochaine commande ?} \\
            }
        }
    }
    & \\ \\

    \theutterance \stepcounter{utterance}  
    & & & \multicolumn{4}{p{0.6\linewidth}}{
        \cellcolor[rgb]{0.9,0.9,0.9}{
            \makecell[{{p{\linewidth}}}]{
                \texttt{\tiny{[GM$\langle$P2]}}
                \texttt{DÉPLACER: P, (4, 4)} \\
            }
        }
    }
     \\ \\

    \theutterance \stepcounter{utterance}  
    & & & \multicolumn{2}{p{0.3\linewidth}}{
        \cellcolor[rgb]{0.9,0.9,0.9}{
            \makecell[{{p{\linewidth}}}]{
                \texttt{\tiny{[GM$|$GM]}}
                \texttt{DÉPLACER: P, (4, 4)} \\
            }
        }
    }
    & & \\ \\

    \theutterance \stepcounter{utterance}  
    & & & \multicolumn{2}{p{0.3\linewidth}}{
        \cellcolor[rgb]{0.9,0.9,0.9}{
            \makecell[{{p{\linewidth}}}]{
                \texttt{\tiny{[GM$|$GM]}}
                \texttt{Pénalité : (4,4) n'est pas vide, mais contient '│'.} \\
\\ 
\texttt{Vous avez collectivement accumulé 1 pénalités sur 8.} \\
\texttt{Veuillez réessayer !} \\
            }
        }
    }
    & & \\ \\

    \theutterance \stepcounter{utterance}  
    & & \multicolumn{4}{p{0.6\linewidth}}{
        \cellcolor[rgb]{0.9,0.9,0.9}{
            \makecell[{{p{\linewidth}}}]{
                \texttt{\tiny{[GM$\rangle$P2]}}
                \texttt{Pénalité : (4,4) n'est pas vide, mais contient '│'.} \\
\\ 
\texttt{Vous avez collectivement accumulé 1 pénalités sur 8.} \\
\texttt{Veuillez réessayer !} \\
            }
        }
    }
    & \\ \\

    \theutterance \stepcounter{utterance}  
    & & & \multicolumn{4}{p{0.6\linewidth}}{
        \cellcolor[rgb]{0.9,0.9,0.9}{
            \makecell[{{p{\linewidth}}}]{
                \texttt{\tiny{[GM$\langle$P2]}}
                \texttt{DIRE: Désolé, (4,4) contient une ligne sur mon plateau, ce n'est pas vide. Je propose (2,4) comme nouvelle cible commune. Je déplace P en (2,4) ce tour. Mes positions restantes: C(4,5), L(4,2). Indiquez{-}moi vos positions et si (2,4) vous convient.} \\
            }
        }
    }
     \\ \\

    \theutterance \stepcounter{utterance}  
    & & & \multicolumn{2}{p{0.3\linewidth}}{
        \cellcolor[rgb]{0.9,0.9,0.9}{
            \makecell[{{p{\linewidth}}}]{
                \texttt{\tiny{[GM$|$GM]}}
                \texttt{DIRE: Désolé, (4,4) contient une ligne sur mon plateau, ce n'est pas vide. Je propose (2,4) comme nouvelle cible commune. Je déplace P en (2,4) ce tour. Mes positions restantes: C(4,5), L(4,2). Indiquez{-}moi vos positions et si (2,4) vous convient.} \\
            }
        }
    }
    & & \\ \\

    \theutterance \stepcounter{utterance}  
    & & \multicolumn{4}{p{0.6\linewidth}}{
        \cellcolor[rgb]{0.9,0.9,0.9}{
            \makecell[{{p{\linewidth}}}]{
                \texttt{\tiny{[P1$\langle$GM]}}
                \texttt{Votre message a été transmis à l'autre joueur.} \\
\\ 
\\ 
\texttt{Vous jouez actuellement le tour 2 sur un maximum de 12.} \\
\texttt{Vous avez collectivement accumulé 1 pénalités sur 8.} \\
\texttt{L'autre joueur a envoyé ce message :} \\
\\ 
\\ 
\texttt{"Désolé, (4,4) contient une ligne sur mon plateau, ce n'est pas vide. Je propose (2,4) comme nouvelle cible commune. Je déplace P en (2,4) ce tour. Mes positions restantes: C(4,5), L(4,2). Indiquez{-}moi vos positions et si (2,4) vous convient."} \\
\\ 
\\ 
\texttt{Quelle est votre prochaine commande ?} \\
            }
        }
    }
    & \\ \\

    \theutterance \stepcounter{utterance}  
    & \multicolumn{4}{p{0.6\linewidth}}{
        \cellcolor[rgb]{0.9,0.9,0.9}{
            \makecell[{{p{\linewidth}}}]{
                \texttt{\tiny{[P1$\rangle$GM]}}
                \texttt{DIRE: Sur mon plateau, (2,4) est libre, ça me convient ! Je déplace mon L en (2,4) ce tour. Mes positions restantes seront : C(1,2) et P(7,7). Pendant que j'occupe (2,4), pouvez{-}vous déplacer C ou L vers une autre case libre ? Proposez{-}moi la prochaine cible. (Note : mon L est déjà en (1,1), je le déplace en (2,4).)} \\
\\ 
\texttt{DÉPLACER: L, (2, 4)} \\
            }
        }
    }
    & & \\ \\

    \theutterance \stepcounter{utterance}  
    & & & \multicolumn{2}{p{0.3\linewidth}}{
        \cellcolor[rgb]{0.9,0.9,0.9}{
            \makecell[{{p{\linewidth}}}]{
                \texttt{\tiny{[GM$|$GM]}}
                \texttt{DIRE: Sur mon plateau, (2,4) est libre, ça me convient ! Je déplace mon L en (2,4) ce tour. Mes positions restantes seront : C(1,2) et P(7,7). Pendant que j'occupe (2,4), pouvez{-}vous déplacer C ou L vers une autre case libre ? Proposez{-}moi la prochaine cible. (Note : mon L est déjà en (1,1), je le déplace en (2,4).)} \\
\\ 
\texttt{DÉPLACER: L, (2, 4)} \\
            }
        }
    }
    & & \\ \\

    \theutterance \stepcounter{utterance}  
    & & & \multicolumn{2}{p{0.3\linewidth}}{
        \cellcolor[rgb]{0.9,0.9,0.9}{
            \makecell[{{p{\linewidth}}}]{
                \texttt{\tiny{[GM$|$GM]}}
                \texttt{Invalid response format: DIRE: Sur mon plateau, (2,4) est libre, ça me convient ! Je déplace mon L en (2,4) ce tour. Mes positions restantes seront : C(1,2) et P(7,7). Pendant que j'occupe (2,4), pouvez{-}vous déplacer C ou L vers une autre case libre ? Proposez{-}moi la prochaine cible. (Note : mon L est déjà en (1,1), je le déplace en (2,4).)} \\
\\ 
\texttt{DÉPLACER: L, (2, 4)} \\
            }
        }
    }
    & & \\ \\

    \theutterance \stepcounter{utterance}  
    & & \multicolumn{4}{p{0.6\linewidth}}{
        \cellcolor[rgb]{0.9,0.9,0.9}{
            \makecell[{{p{\linewidth}}}]{
                \texttt{\tiny{[P1$\langle$GM]}}
                \texttt{Pénalité : Votre message contient plus d'une commande !} \\
\\ 
\texttt{Assurez{-}vous que votre réponse ne contienne que `DIRE: <MESSAGE>` ou `DÉPLACER: <OBJET>, (<X>, <Y>)`, et rien d'autre !} \\
\\ 
\texttt{Vous avez collectivement accumulé 2 pénalités sur 8.} \\
\texttt{Veuillez réessayer !} \\
            }
        }
    }
    & \\ \\

    \theutterance \stepcounter{utterance}  
    & \multicolumn{4}{p{0.6\linewidth}}{
        \cellcolor[rgb]{0.9,0.9,0.9}{
            \makecell[{{p{\linewidth}}}]{
                \texttt{\tiny{[P1$\rangle$GM]}}
                \texttt{DÉPLACER: L, (2, 4)} \\
            }
        }
    }
    & & \\ \\

    \theutterance \stepcounter{utterance}  
    & & & \multicolumn{2}{p{0.3\linewidth}}{
        \cellcolor[rgb]{0.9,0.9,0.9}{
            \makecell[{{p{\linewidth}}}]{
                \texttt{\tiny{[GM$|$GM]}}
                \texttt{DÉPLACER: L, (2, 4)} \\
            }
        }
    }
    & & \\ \\

    \theutterance \stepcounter{utterance}  
    & & & \multicolumn{2}{p{0.3\linewidth}}{
        \cellcolor[rgb]{0.9,0.9,0.9}{
            \makecell[{{p{\linewidth}}}]{
                \texttt{\tiny{[GM$|$GM]}}
                \texttt{Pénalité : (2,4) n'est pas vide, mais contient '│'.} \\
\\ 
\texttt{Vous avez collectivement accumulé 3 pénalités sur 8.} \\
\texttt{Veuillez réessayer !} \\
            }
        }
    }
    & & \\ \\

    \theutterance \stepcounter{utterance}  
    & & \multicolumn{4}{p{0.6\linewidth}}{
        \cellcolor[rgb]{0.9,0.9,0.9}{
            \makecell[{{p{\linewidth}}}]{
                \texttt{\tiny{[P1$\langle$GM]}}
                \texttt{Pénalité : (2,4) n'est pas vide, mais contient '│'.} \\
\\ 
\texttt{Vous avez collectivement accumulé 3 pénalités sur 8.} \\
\texttt{Veuillez réessayer !} \\
            }
        }
    }
    & \\ \\

    \theutterance \stepcounter{utterance}  
    & \multicolumn{4}{p{0.6\linewidth}}{
        \cellcolor[rgb]{0.9,0.9,0.9}{
            \makecell[{{p{\linewidth}}}]{
                \texttt{\tiny{[P1$\rangle$GM]}}
                \texttt{DÉPLACER: L, (4, 4)} \\
            }
        }
    }
    & & \\ \\

    \theutterance \stepcounter{utterance}  
    & & & \multicolumn{2}{p{0.3\linewidth}}{
        \cellcolor[rgb]{0.9,0.9,0.9}{
            \makecell[{{p{\linewidth}}}]{
                \texttt{\tiny{[GM$|$GM]}}
                \texttt{DÉPLACER: L, (4, 4)} \\
            }
        }
    }
    & & \\ \\

    \theutterance \stepcounter{utterance}  
    & & & \multicolumn{2}{p{0.3\linewidth}}{
        \cellcolor[rgb]{0.9,0.9,0.9}{
            \makecell[{{p{\linewidth}}}]{
                \texttt{\tiny{[GM$|$GM]}}
                \texttt{Pénalité : (4,4) n'est pas vide, mais contient '│'.} \\
\\ 
\texttt{Vous avez collectivement accumulé 4 pénalités sur 8.} \\
\texttt{Veuillez réessayer !} \\
            }
        }
    }
    & & \\ \\

    \theutterance \stepcounter{utterance}  
    & & \multicolumn{4}{p{0.6\linewidth}}{
        \cellcolor[rgb]{0.9,0.9,0.9}{
            \makecell[{{p{\linewidth}}}]{
                \texttt{\tiny{[P1$\langle$GM]}}
                \texttt{Pénalité : (4,4) n'est pas vide, mais contient '│'.} \\
\\ 
\texttt{Vous avez collectivement accumulé 4 pénalités sur 8.} \\
\texttt{Veuillez réessayer !} \\
            }
        }
    }
    & \\ \\

    \theutterance \stepcounter{utterance}  
    & \multicolumn{4}{p{0.6\linewidth}}{
        \cellcolor[rgb]{0.9,0.9,0.9}{
            \makecell[{{p{\linewidth}}}]{
                \texttt{\tiny{[P1$\rangle$GM]}}
                \texttt{DIRE: Désolé pour les erreurs, je vois maintenant que (2,4) et (4,4) contiennent des lignes sur mon plateau. Je propose (3,4) comme cible commune, qui est vide chez moi. Mes positions actuelles : L(1,1), C(1,2), P(7,7). Pouvez{-}vous confirmer si (3,4) est libre chez vous et déplacer un objet là{-}bas ?} \\
            }
        }
    }
    & & \\ \\

    \theutterance \stepcounter{utterance}  
    & & & \multicolumn{2}{p{0.3\linewidth}}{
        \cellcolor[rgb]{0.9,0.9,0.9}{
            \makecell[{{p{\linewidth}}}]{
                \texttt{\tiny{[GM$|$GM]}}
                \texttt{DIRE: Désolé pour les erreurs, je vois maintenant que (2,4) et (4,4) contiennent des lignes sur mon plateau. Je propose (3,4) comme cible commune, qui est vide chez moi. Mes positions actuelles : L(1,1), C(1,2), P(7,7). Pouvez{-}vous confirmer si (3,4) est libre chez vous et déplacer un objet là{-}bas ?} \\
            }
        }
    }
    & & \\ \\

    \theutterance \stepcounter{utterance}  
    & & \multicolumn{4}{p{0.6\linewidth}}{
        \cellcolor[rgb]{0.9,0.9,0.9}{
            \makecell[{{p{\linewidth}}}]{
                \texttt{\tiny{[GM$\rangle$P2]}}
                \texttt{Votre message a été transmis à l'autre joueur.} \\
\\ 
\\ 
\texttt{Vous jouez actuellement le tour 3 sur un maximum de 12.} \\
\texttt{Vous avez collectivement accumulé 4 pénalités sur 8.} \\
\texttt{L'autre joueur a envoyé ce message :} \\
\\ 
\\ 
\texttt{"Désolé pour les erreurs, je vois maintenant que (2,4) et (4,4) contiennent des lignes sur mon plateau. Je propose (3,4) comme cible commune, qui est vide chez moi. Mes positions actuelles : L(1,1), C(1,2), P(7,7). Pouvez{-}vous confirmer si (3,4) est libre chez vous et déplacer un objet là{-}bas ?"} \\
\\ 
\\ 
\texttt{Quelle est votre prochaine commande ?} \\
            }
        }
    }
    & \\ \\

    \theutterance \stepcounter{utterance}  
    & & & \multicolumn{4}{p{0.6\linewidth}}{
        \cellcolor[rgb]{0.9,0.9,0.9}{
            \makecell[{{p{\linewidth}}}]{
                \texttt{\tiny{[GM$\langle$P2]}}
                \texttt{DÉPLACER: P, (3, 4)} \\
            }
        }
    }
     \\ \\

    \theutterance \stepcounter{utterance}  
    & & & \multicolumn{2}{p{0.3\linewidth}}{
        \cellcolor[rgb]{0.9,0.9,0.9}{
            \makecell[{{p{\linewidth}}}]{
                \texttt{\tiny{[GM$|$GM]}}
                \texttt{DÉPLACER: P, (3, 4)} \\
            }
        }
    }
    & & \\ \\

    \theutterance \stepcounter{utterance}  
    & & & \multicolumn{2}{p{0.3\linewidth}}{
        \cellcolor[rgb]{0.9,0.9,0.9}{
            \makecell[{{p{\linewidth}}}]{
                \texttt{\tiny{[GM$|$GM]}}
                \texttt{P moved to (3, 4)} \\
            }
        }
    }
    & & \\ \\

    \theutterance \stepcounter{utterance}  
    & & \multicolumn{4}{p{0.6\linewidth}}{
        \cellcolor[rgb]{0.9,0.9,0.9}{
            \makecell[{{p{\linewidth}}}]{
                \texttt{\tiny{[P1$\langle$GM]}}
                \texttt{Votre message a été transmis à l'autre joueur.} \\
\\ 
\\ 
\texttt{Vous jouez actuellement le tour 3 sur un maximum de 12.} \\
\texttt{Vous avez collectivement accumulé 4 pénalités sur 8.} \\
\texttt{L'autre joueur a déplacé un objet sur sa grille.} \\
\\ 
\\ 
\texttt{Quelle est votre prochaine commande ?} \\
            }
        }
    }
    & \\ \\

    \theutterance \stepcounter{utterance}  
    & \multicolumn{4}{p{0.6\linewidth}}{
        \cellcolor[rgb]{0.9,0.9,0.9}{
            \makecell[{{p{\linewidth}}}]{
                \texttt{\tiny{[P1$\rangle$GM]}}
                \texttt{DIRE: Je n'ai pas eu de réponse sur (3,4). Pouvez{-}vous me dire si (3,4) est libre chez vous ? Je voudrais aussi connaître vos positions actuelles exactes pour synchroniser nos déplacements. Mes positions: L(1,1), C(1,2), P(7,7).} \\
            }
        }
    }
    & & \\ \\

    \theutterance \stepcounter{utterance}  
    & & & \multicolumn{2}{p{0.3\linewidth}}{
        \cellcolor[rgb]{0.9,0.9,0.9}{
            \makecell[{{p{\linewidth}}}]{
                \texttt{\tiny{[GM$|$GM]}}
                \texttt{DIRE: Je n'ai pas eu de réponse sur (3,4). Pouvez{-}vous me dire si (3,4) est libre chez vous ? Je voudrais aussi connaître vos positions actuelles exactes pour synchroniser nos déplacements. Mes positions: L(1,1), C(1,2), P(7,7).} \\
            }
        }
    }
    & & \\ \\

    \theutterance \stepcounter{utterance}  
    & & \multicolumn{4}{p{0.6\linewidth}}{
        \cellcolor[rgb]{0.9,0.9,0.9}{
            \makecell[{{p{\linewidth}}}]{
                \texttt{\tiny{[GM$\rangle$P2]}}
                \texttt{Déplacement de 'P' vers (3,4) réussi. Votre grille mise à jour ressemble à ceci :} \\
\texttt{```} \\
\texttt{ 1234567} \\
\texttt{╔═╤═╤═══╗} \\
\texttt{║◌│◌│◌◌◌║ 1} \\
\texttt{║◌├─┤◌L◌║ 2} \\
\texttt{╟─┤◌├─┬─╢ 3} \\
\texttt{║◌│P│◌│◌║ 4} \\
\texttt{╟─┼─┘C│◌║ 5} \\
\texttt{║◌│◌◌┌┴─╢ 6} \\
\texttt{║◌│◌◌│◌◌║ 7} \\
\texttt{╚═╧══╧══╝} \\
\texttt{```} \\
\texttt{Vous jouez actuellement le tour 4 sur un maximum de 12.} \\
\texttt{Vous avez collectivement accumulé 4 pénalités sur 8.} \\
\texttt{L'autre joueur a envoyé ce message :} \\
\\ 
\\ 
\texttt{"Je n'ai pas eu de réponse sur (3,4). Pouvez{-}vous me dire si (3,4) est libre chez vous ? Je voudrais aussi connaître vos positions actuelles exactes pour synchroniser nos déplacements. Mes positions: L(1,1), C(1,2), P(7,7)."} \\
\\ 
\\ 
\texttt{Quelle est votre prochaine commande ?} \\
            }
        }
    }
    & \\ \\

    \theutterance \stepcounter{utterance}  
    & & & \multicolumn{4}{p{0.6\linewidth}}{
        \cellcolor[rgb]{0.9,0.9,0.9}{
            \makecell[{{p{\linewidth}}}]{
                \texttt{\tiny{[GM$\langle$P2]}}
                \texttt{DIRE: P est maintenant en (3,4), c'était libre. Mes positions actuelles: P(3,4), C(4,5), L(4,2). Je vous propose de placer L en (1,4) et C en (2,4) comme cibles communes. Pouvez{-}vous déplacer votre L vers (1,4) et votre C vers (2,4) ?} \\
            }
        }
    }
     \\ \\

    \theutterance \stepcounter{utterance}  
    & & & \multicolumn{2}{p{0.3\linewidth}}{
        \cellcolor[rgb]{0.9,0.9,0.9}{
            \makecell[{{p{\linewidth}}}]{
                \texttt{\tiny{[GM$|$GM]}}
                \texttt{DIRE: P est maintenant en (3,4), c'était libre. Mes positions actuelles: P(3,4), C(4,5), L(4,2). Je vous propose de placer L en (1,4) et C en (2,4) comme cibles communes. Pouvez{-}vous déplacer votre L vers (1,4) et votre C vers (2,4) ?} \\
            }
        }
    }
    & & \\ \\

    \theutterance \stepcounter{utterance}  
    & & \multicolumn{4}{p{0.6\linewidth}}{
        \cellcolor[rgb]{0.9,0.9,0.9}{
            \makecell[{{p{\linewidth}}}]{
                \texttt{\tiny{[P1$\langle$GM]}}
                \texttt{Votre message a été transmis à l'autre joueur.} \\
\\ 
\\ 
\texttt{Vous jouez actuellement le tour 4 sur un maximum de 12.} \\
\texttt{Vous avez collectivement accumulé 4 pénalités sur 8.} \\
\texttt{L'autre joueur a envoyé ce message :} \\
\\ 
\\ 
\texttt{"P est maintenant en (3,4), c'était libre. Mes positions actuelles: P(3,4), C(4,5), L(4,2). Je vous propose de placer L en (1,4) et C en (2,4) comme cibles communes. Pouvez{-}vous déplacer votre L vers (1,4) et votre C vers (2,4) ?"} \\
\\ 
\\ 
\texttt{Quelle est votre prochaine commande ?} \\
            }
        }
    }
    & \\ \\

    \theutterance \stepcounter{utterance}  
    & \multicolumn{4}{p{0.6\linewidth}}{
        \cellcolor[rgb]{0.9,0.9,0.9}{
            \makecell[{{p{\linewidth}}}]{
                \texttt{\tiny{[P1$\rangle$GM]}}
                \texttt{DIRE: Malheureusement, (1,4) n'est pas vide sur mon plateau (contient '│') non plus. Pour synchroniser, je vous propose de me donner une liste de plusieurs coordonnées libres chez vous, et je vérifierai lesquelles sont aussi libres chez moi. Mes cases qui semblent vides incluent notamment (5,4), (6,4), (7,4), (5,3), (6,5), (7,5), (6,6), (7,6). Pouvez{-}vous me dire lesquelles de ces coordonnées sont libres sur votre plateau ?} \\
            }
        }
    }
    & & \\ \\

    \theutterance \stepcounter{utterance}  
    & & & \multicolumn{2}{p{0.3\linewidth}}{
        \cellcolor[rgb]{0.9,0.9,0.9}{
            \makecell[{{p{\linewidth}}}]{
                \texttt{\tiny{[GM$|$GM]}}
                \texttt{DIRE: Malheureusement, (1,4) n'est pas vide sur mon plateau (contient '│') non plus. Pour synchroniser, je vous propose de me donner une liste de plusieurs coordonnées libres chez vous, et je vérifierai lesquelles sont aussi libres chez moi. Mes cases qui semblent vides incluent notamment (5,4), (6,4), (7,4), (5,3), (6,5), (7,5), (6,6), (7,6). Pouvez{-}vous me dire lesquelles de ces coordonnées sont libres sur votre plateau ?} \\
            }
        }
    }
    & & \\ \\

    \theutterance \stepcounter{utterance}  
    & & \multicolumn{4}{p{0.6\linewidth}}{
        \cellcolor[rgb]{0.9,0.9,0.9}{
            \makecell[{{p{\linewidth}}}]{
                \texttt{\tiny{[GM$\rangle$P2]}}
                \texttt{Votre message a été transmis à l'autre joueur.} \\
\\ 
\\ 
\texttt{Vous jouez actuellement le tour 5 sur un maximum de 12.} \\
\texttt{Vous avez collectivement accumulé 4 pénalités sur 8.} \\
\texttt{L'autre joueur a envoyé ce message :} \\
\\ 
\\ 
\texttt{"Malheureusement, (1,4) n'est pas vide sur mon plateau (contient '│') non plus. Pour synchroniser, je vous propose de me donner une liste de plusieurs coordonnées libres chez vous, et je vérifierai lesquelles sont aussi libres chez moi. Mes cases qui semblent vides incluent notamment (5,4), (6,4), (7,4), (5,3), (6,5), (7,5), (6,6), (7,6). Pouvez{-}vous me dire lesquelles de ces coordonnées sont libres sur votre plateau ?"} \\
\\ 
\\ 
\texttt{Quelle est votre prochaine commande ?} \\
            }
        }
    }
    & \\ \\

    \theutterance \stepcounter{utterance}  
    & & & \multicolumn{4}{p{0.6\linewidth}}{
        \cellcolor[rgb]{0.9,0.9,0.9}{
            \makecell[{{p{\linewidth}}}]{
                \texttt{\tiny{[GM$\langle$P2]}}
                \texttt{DIRE: Voici les coordonnées libres sur mon plateau parmi votre liste: (5,4), (6,4), (7,4), (5,3), (6,5), (7,5), (6,6), (7,6) sont toutes libres. Je propose comme cibles finales: P en (3,4), L en (5,3), C en (6,5). Pouvez{-}vous confirmer que (5,3) et (6,5) sont libres chez vous? Je vais déplacer L vers (5,3) ce tour.} \\
            }
        }
    }
     \\ \\

    \theutterance \stepcounter{utterance}  
    & & & \multicolumn{2}{p{0.3\linewidth}}{
        \cellcolor[rgb]{0.9,0.9,0.9}{
            \makecell[{{p{\linewidth}}}]{
                \texttt{\tiny{[GM$|$GM]}}
                \texttt{DIRE: Voici les coordonnées libres sur mon plateau parmi votre liste: (5,4), (6,4), (7,4), (5,3), (6,5), (7,5), (6,6), (7,6) sont toutes libres. Je propose comme cibles finales: P en (3,4), L en (5,3), C en (6,5). Pouvez{-}vous confirmer que (5,3) et (6,5) sont libres chez vous? Je vais déplacer L vers (5,3) ce tour.} \\
            }
        }
    }
    & & \\ \\

    \theutterance \stepcounter{utterance}  
    & & \multicolumn{4}{p{0.6\linewidth}}{
        \cellcolor[rgb]{0.9,0.9,0.9}{
            \makecell[{{p{\linewidth}}}]{
                \texttt{\tiny{[P1$\langle$GM]}}
                \texttt{Votre message a été transmis à l'autre joueur.} \\
\\ 
\\ 
\texttt{Vous jouez actuellement le tour 5 sur un maximum de 12.} \\
\texttt{Vous avez collectivement accumulé 4 pénalités sur 8.} \\
\texttt{L'autre joueur a envoyé ce message :} \\
\\ 
\\ 
\texttt{"Voici les coordonnées libres sur mon plateau parmi votre liste: (5,4), (6,4), (7,4), (5,3), (6,5), (7,5), (6,6), (7,6) sont toutes libres. Je propose comme cibles finales: P en (3,4), L en (5,3), C en (6,5). Pouvez{-}vous confirmer que (5,3) et (6,5) sont libres chez vous? Je vais déplacer L vers (5,3) ce tour."} \\
\\ 
\\ 
\texttt{Quelle est votre prochaine commande ?} \\
            }
        }
    }
    & \\ \\

    \theutterance \stepcounter{utterance}  
    & \multicolumn{4}{p{0.6\linewidth}}{
        \cellcolor[rgb]{0.9,0.9,0.9}{
            \makecell[{{p{\linewidth}}}]{
                \texttt{\tiny{[P1$\rangle$GM]}}
                \texttt{DIRE: (5,3) et (6,5) sont libres sur mon plateau. C'est parfait ! Je déplace mon L en (5,3) ce tour. La prochaine cible sera C en (6,5). Faites de même avec votre C le prochain tour.} \\
\\ 
\texttt{DÉPLACER: L, (5, 3)} \\
            }
        }
    }
    & & \\ \\

    \theutterance \stepcounter{utterance}  
    & & & \multicolumn{2}{p{0.3\linewidth}}{
        \cellcolor[rgb]{0.9,0.9,0.9}{
            \makecell[{{p{\linewidth}}}]{
                \texttt{\tiny{[GM$|$GM]}}
                \texttt{DIRE: (5,3) et (6,5) sont libres sur mon plateau. C'est parfait ! Je déplace mon L en (5,3) ce tour. La prochaine cible sera C en (6,5). Faites de même avec votre C le prochain tour.} \\
\\ 
\texttt{DÉPLACER: L, (5, 3)} \\
            }
        }
    }
    & & \\ \\

    \theutterance \stepcounter{utterance}  
    & & & \multicolumn{2}{p{0.3\linewidth}}{
        \cellcolor[rgb]{0.9,0.9,0.9}{
            \makecell[{{p{\linewidth}}}]{
                \texttt{\tiny{[GM$|$GM]}}
                \texttt{Invalid response format: DIRE: (5,3) et (6,5) sont libres sur mon plateau. C'est parfait ! Je déplace mon L en (5,3) ce tour. La prochaine cible sera C en (6,5). Faites de même avec votre C le prochain tour.} \\
\\ 
\texttt{DÉPLACER: L, (5, 3)} \\
            }
        }
    }
    & & \\ \\

    \theutterance \stepcounter{utterance}  
    & & \multicolumn{4}{p{0.6\linewidth}}{
        \cellcolor[rgb]{0.9,0.9,0.9}{
            \makecell[{{p{\linewidth}}}]{
                \texttt{\tiny{[P1$\langle$GM]}}
                \texttt{Pénalité : Votre message contient plus d'une commande !} \\
\\ 
\texttt{Assurez{-}vous que votre réponse ne contienne que `DIRE: <MESSAGE>` ou `DÉPLACER: <OBJET>, (<X>, <Y>)`, et rien d'autre !} \\
\\ 
\texttt{Vous avez collectivement accumulé 5 pénalités sur 8.} \\
\texttt{Veuillez réessayer !} \\
            }
        }
    }
    & \\ \\

    \theutterance \stepcounter{utterance}  
    & \multicolumn{4}{p{0.6\linewidth}}{
        \cellcolor[rgb]{0.9,0.9,0.9}{
            \makecell[{{p{\linewidth}}}]{
                \texttt{\tiny{[P1$\rangle$GM]}}
                \texttt{DÉPLACER: L, (5, 3)} \\
            }
        }
    }
    & & \\ \\

    \theutterance \stepcounter{utterance}  
    & & & \multicolumn{2}{p{0.3\linewidth}}{
        \cellcolor[rgb]{0.9,0.9,0.9}{
            \makecell[{{p{\linewidth}}}]{
                \texttt{\tiny{[GM$|$GM]}}
                \texttt{DÉPLACER: L, (5, 3)} \\
            }
        }
    }
    & & \\ \\

    \theutterance \stepcounter{utterance}  
    & & & \multicolumn{2}{p{0.3\linewidth}}{
        \cellcolor[rgb]{0.9,0.9,0.9}{
            \makecell[{{p{\linewidth}}}]{
                \texttt{\tiny{[GM$|$GM]}}
                \texttt{Pénalité : (5,3) n'est pas vide, mais contient '─'.} \\
\\ 
\texttt{Vous avez collectivement accumulé 6 pénalités sur 8.} \\
\texttt{Veuillez réessayer !} \\
            }
        }
    }
    & & \\ \\

    \theutterance \stepcounter{utterance}  
    & & \multicolumn{4}{p{0.6\linewidth}}{
        \cellcolor[rgb]{0.9,0.9,0.9}{
            \makecell[{{p{\linewidth}}}]{
                \texttt{\tiny{[P1$\langle$GM]}}
                \texttt{Pénalité : (5,3) n'est pas vide, mais contient '─'.} \\
\\ 
\texttt{Vous avez collectivement accumulé 6 pénalités sur 8.} \\
\texttt{Veuillez réessayer !} \\
            }
        }
    }
    & \\ \\

    \theutterance \stepcounter{utterance}  
    & \multicolumn{4}{p{0.6\linewidth}}{
        \cellcolor[rgb]{0.9,0.9,0.9}{
            \makecell[{{p{\linewidth}}}]{
                \texttt{\tiny{[P1$\rangle$GM]}}
                \texttt{DIRE: Désolé, (5,3) n'est pas vide sur mon plateau (contient '─'). Parmi mes cases réellement vides, j'ai vérifié : (5,4), (6,4), (6,5), (6,6), (7,6) semblent vides. Je vous propose comme cibles : P en (3,4), L en (5,4), C en (6,5). Pouvez{-}vous déplacer votre L vers (5,4) et confirmer ?} \\
            }
        }
    }
    & & \\ \\

    \theutterance \stepcounter{utterance}  
    & & & \multicolumn{2}{p{0.3\linewidth}}{
        \cellcolor[rgb]{0.9,0.9,0.9}{
            \makecell[{{p{\linewidth}}}]{
                \texttt{\tiny{[GM$|$GM]}}
                \texttt{DIRE: Désolé, (5,3) n'est pas vide sur mon plateau (contient '─'). Parmi mes cases réellement vides, j'ai vérifié : (5,4), (6,4), (6,5), (6,6), (7,6) semblent vides. Je vous propose comme cibles : P en (3,4), L en (5,4), C en (6,5). Pouvez{-}vous déplacer votre L vers (5,4) et confirmer ?} \\
            }
        }
    }
    & & \\ \\

    \theutterance \stepcounter{utterance}  
    & & \multicolumn{4}{p{0.6\linewidth}}{
        \cellcolor[rgb]{0.9,0.9,0.9}{
            \makecell[{{p{\linewidth}}}]{
                \texttt{\tiny{[GM$\rangle$P2]}}
                \texttt{Votre message a été transmis à l'autre joueur.} \\
\\ 
\\ 
\texttt{Vous jouez actuellement le tour 6 sur un maximum de 12.} \\
\texttt{Vous avez collectivement accumulé 6 pénalités sur 8.} \\
\texttt{L'autre joueur a envoyé ce message :} \\
\\ 
\\ 
\texttt{"Désolé, (5,3) n'est pas vide sur mon plateau (contient '─'). Parmi mes cases réellement vides, j'ai vérifié : (5,4), (6,4), (6,5), (6,6), (7,6) semblent vides. Je vous propose comme cibles : P en (3,4), L en (5,4), C en (6,5). Pouvez{-}vous déplacer votre L vers (5,4) et confirmer ?"} \\
\\ 
\\ 
\texttt{Quelle est votre prochaine commande ?} \\
            }
        }
    }
    & \\ \\

    \theutterance \stepcounter{utterance}  
    & & & \multicolumn{4}{p{0.6\linewidth}}{
        \cellcolor[rgb]{0.9,0.9,0.9}{
            \makecell[{{p{\linewidth}}}]{
                \texttt{\tiny{[GM$\langle$P2]}}
                \texttt{DÉPLACER: L, (5, 4)} \\
            }
        }
    }
     \\ \\

    \theutterance \stepcounter{utterance}  
    & & & \multicolumn{2}{p{0.3\linewidth}}{
        \cellcolor[rgb]{0.9,0.9,0.9}{
            \makecell[{{p{\linewidth}}}]{
                \texttt{\tiny{[GM$|$GM]}}
                \texttt{DÉPLACER: L, (5, 4)} \\
            }
        }
    }
    & & \\ \\

    \theutterance \stepcounter{utterance}  
    & & & \multicolumn{2}{p{0.3\linewidth}}{
        \cellcolor[rgb]{0.9,0.9,0.9}{
            \makecell[{{p{\linewidth}}}]{
                \texttt{\tiny{[GM$|$GM]}}
                \texttt{L moved to (5, 4)} \\
            }
        }
    }
    & & \\ \\

    \theutterance \stepcounter{utterance}  
    & & \multicolumn{4}{p{0.6\linewidth}}{
        \cellcolor[rgb]{0.9,0.9,0.9}{
            \makecell[{{p{\linewidth}}}]{
                \texttt{\tiny{[P1$\langle$GM]}}
                \texttt{Votre message a été transmis à l'autre joueur.} \\
\\ 
\\ 
\texttt{Vous jouez actuellement le tour 6 sur un maximum de 12.} \\
\texttt{Vous avez collectivement accumulé 6 pénalités sur 8.} \\
\texttt{L'autre joueur a déplacé un objet sur sa grille.} \\
\\ 
\\ 
\texttt{Quelle est votre prochaine commande ?} \\
            }
        }
    }
    & \\ \\

    \theutterance \stepcounter{utterance}  
    & \multicolumn{4}{p{0.6\linewidth}}{
        \cellcolor[rgb]{0.9,0.9,0.9}{
            \makecell[{{p{\linewidth}}}]{
                \texttt{\tiny{[P1$\rangle$GM]}}
                \texttt{DÉPLACER: L, (5, 4)} \\
            }
        }
    }
    & & \\ \\

    \theutterance \stepcounter{utterance}  
    & & & \multicolumn{2}{p{0.3\linewidth}}{
        \cellcolor[rgb]{0.9,0.9,0.9}{
            \makecell[{{p{\linewidth}}}]{
                \texttt{\tiny{[GM$|$GM]}}
                \texttt{DÉPLACER: L, (5, 4)} \\
            }
        }
    }
    & & \\ \\

    \theutterance \stepcounter{utterance}  
    & & & \multicolumn{2}{p{0.3\linewidth}}{
        \cellcolor[rgb]{0.9,0.9,0.9}{
            \makecell[{{p{\linewidth}}}]{
                \texttt{\tiny{[GM$|$GM]}}
                \texttt{L moved to (5, 4)} \\
            }
        }
    }
    & & \\ \\

    \theutterance \stepcounter{utterance}  
    & & \multicolumn{4}{p{0.6\linewidth}}{
        \cellcolor[rgb]{0.9,0.9,0.9}{
            \makecell[{{p{\linewidth}}}]{
                \texttt{\tiny{[GM$\rangle$P2]}}
                \texttt{Déplacement de 'L' vers (5,4) réussi. Votre grille mise à jour ressemble à ceci :} \\
\texttt{```} \\
\texttt{ 1234567} \\
\texttt{╔═╤═╤═══╗} \\
\texttt{║◌│◌│◌◌◌║ 1} \\
\texttt{║◌├─┤◌◌◌║ 2} \\
\texttt{╟─┤◌├─┬─╢ 3} \\
\texttt{║◌│P│L│◌║ 4} \\
\texttt{╟─┼─┘C│◌║ 5} \\
\texttt{║◌│◌◌┌┴─╢ 6} \\
\texttt{║◌│◌◌│◌◌║ 7} \\
\texttt{╚═╧══╧══╝} \\
\texttt{```} \\
\texttt{Vous jouez actuellement le tour 7 sur un maximum de 12.} \\
\texttt{Vous avez collectivement accumulé 6 pénalités sur 8.} \\
\texttt{L'autre joueur a déplacé un objet sur sa grille.} \\
\\ 
\\ 
\texttt{Quelle est votre prochaine commande ?} \\
            }
        }
    }
    & \\ \\

    \theutterance \stepcounter{utterance}  
    & & & \multicolumn{4}{p{0.6\linewidth}}{
        \cellcolor[rgb]{0.9,0.9,0.9}{
            \makecell[{{p{\linewidth}}}]{
                \texttt{\tiny{[GM$\langle$P2]}}
                \texttt{DIRE: L est maintenant en (5,4). Mes positions actuelles: P(3,4), L(5,4), C(4,5). Je vais déplacer C vers (6,5) ce tour. Pouvez{-}vous confirmer que vous placez aussi vos objets aux mêmes coordonnées: P(3,4), L(5,4), C(6,5)?} \\
            }
        }
    }
     \\ \\

    \theutterance \stepcounter{utterance}  
    & & & \multicolumn{2}{p{0.3\linewidth}}{
        \cellcolor[rgb]{0.9,0.9,0.9}{
            \makecell[{{p{\linewidth}}}]{
                \texttt{\tiny{[GM$|$GM]}}
                \texttt{DIRE: L est maintenant en (5,4). Mes positions actuelles: P(3,4), L(5,4), C(4,5). Je vais déplacer C vers (6,5) ce tour. Pouvez{-}vous confirmer que vous placez aussi vos objets aux mêmes coordonnées: P(3,4), L(5,4), C(6,5)?} \\
            }
        }
    }
    & & \\ \\

    \theutterance \stepcounter{utterance}  
    & & \multicolumn{4}{p{0.6\linewidth}}{
        \cellcolor[rgb]{0.9,0.9,0.9}{
            \makecell[{{p{\linewidth}}}]{
                \texttt{\tiny{[P1$\langle$GM]}}
                \texttt{Déplacement de 'L' vers (5,4) réussi. Votre grille mise à jour ressemble à ceci :} \\
\texttt{```} \\
\texttt{ 1234567} \\
\texttt{╔═╤═╤═══╗} \\
\texttt{║◌│◌│◌◌◌║ 1} \\
\texttt{║C├─┤◌◌◌║ 2} \\
\texttt{╟─┤◌├─┬─╢ 3} \\
\texttt{║◌│◌│L│◌║ 4} \\
\texttt{╟─┼─┘◌│◌║ 5} \\
\texttt{║◌│◌◌┌┴─╢ 6} \\
\texttt{║◌│◌◌│◌P║ 7} \\
\texttt{╚═╧══╧══╝} \\
\texttt{```} \\
\texttt{Vous jouez actuellement le tour 7 sur un maximum de 12.} \\
\texttt{Vous avez collectivement accumulé 6 pénalités sur 8.} \\
\texttt{L'autre joueur a envoyé ce message :} \\
\\ 
\\ 
\texttt{"L est maintenant en (5,4). Mes positions actuelles: P(3,4), L(5,4), C(4,5). Je vais déplacer C vers (6,5) ce tour. Pouvez{-}vous confirmer que vous placez aussi vos objets aux mêmes coordonnées: P(3,4), L(5,4), C(6,5)?"} \\
\\ 
\\ 
\texttt{Quelle est votre prochaine commande ?} \\
            }
        }
    }
    & \\ \\

    \theutterance \stepcounter{utterance}  
    & \multicolumn{4}{p{0.6\linewidth}}{
        \cellcolor[rgb]{0.9,0.9,0.9}{
            \makecell[{{p{\linewidth}}}]{
                \texttt{\tiny{[P1$\rangle$GM]}}
                \texttt{DIRE: Oui, je confirme la configuration finale : P(3,4), L(5,4), C(6,5). Mon L est déjà en (5,4). Je vais déplacer mon C vers (6,5) ce tour, puis P vers (3,4) au tour suivant. Faites de même pour votre P et C.} \\
            }
        }
    }
    & & \\ \\

    \theutterance \stepcounter{utterance}  
    & & & \multicolumn{2}{p{0.3\linewidth}}{
        \cellcolor[rgb]{0.9,0.9,0.9}{
            \makecell[{{p{\linewidth}}}]{
                \texttt{\tiny{[GM$|$GM]}}
                \texttt{DIRE: Oui, je confirme la configuration finale : P(3,4), L(5,4), C(6,5). Mon L est déjà en (5,4). Je vais déplacer mon C vers (6,5) ce tour, puis P vers (3,4) au tour suivant. Faites de même pour votre P et C.} \\
            }
        }
    }
    & & \\ \\

    \theutterance \stepcounter{utterance}  
    & & \multicolumn{4}{p{0.6\linewidth}}{
        \cellcolor[rgb]{0.9,0.9,0.9}{
            \makecell[{{p{\linewidth}}}]{
                \texttt{\tiny{[GM$\rangle$P2]}}
                \texttt{Votre message a été transmis à l'autre joueur.} \\
\\ 
\\ 
\texttt{Vous jouez actuellement le tour 8 sur un maximum de 12.} \\
\texttt{Vous avez collectivement accumulé 6 pénalités sur 8.} \\
\texttt{L'autre joueur a envoyé ce message :} \\
\\ 
\\ 
\texttt{"Oui, je confirme la configuration finale : P(3,4), L(5,4), C(6,5). Mon L est déjà en (5,4). Je vais déplacer mon C vers (6,5) ce tour, puis P vers (3,4) au tour suivant. Faites de même pour votre P et C."} \\
\\ 
\\ 
\texttt{Quelle est votre prochaine commande ?} \\
            }
        }
    }
    & \\ \\

    \theutterance \stepcounter{utterance}  
    & & & \multicolumn{4}{p{0.6\linewidth}}{
        \cellcolor[rgb]{0.9,0.9,0.9}{
            \makecell[{{p{\linewidth}}}]{
                \texttt{\tiny{[GM$\langle$P2]}}
                \texttt{DÉPLACER: C, (6, 5)} \\
            }
        }
    }
     \\ \\

    \theutterance \stepcounter{utterance}  
    & & & \multicolumn{2}{p{0.3\linewidth}}{
        \cellcolor[rgb]{0.9,0.9,0.9}{
            \makecell[{{p{\linewidth}}}]{
                \texttt{\tiny{[GM$|$GM]}}
                \texttt{DÉPLACER: C, (6, 5)} \\
            }
        }
    }
    & & \\ \\

    \theutterance \stepcounter{utterance}  
    & & & \multicolumn{2}{p{0.3\linewidth}}{
        \cellcolor[rgb]{0.9,0.9,0.9}{
            \makecell[{{p{\linewidth}}}]{
                \texttt{\tiny{[GM$|$GM]}}
                \texttt{Pénalité : (6,5) n'est pas vide, mais contient '│'.} \\
\\ 
\texttt{Vous avez collectivement accumulé 7 pénalités sur 8.} \\
\texttt{Veuillez réessayer !} \\
            }
        }
    }
    & & \\ \\

    \theutterance \stepcounter{utterance}  
    & & \multicolumn{4}{p{0.6\linewidth}}{
        \cellcolor[rgb]{0.9,0.9,0.9}{
            \makecell[{{p{\linewidth}}}]{
                \texttt{\tiny{[GM$\rangle$P2]}}
                \texttt{Pénalité : (6,5) n'est pas vide, mais contient '│'.} \\
\\ 
\texttt{Vous avez collectivement accumulé 7 pénalités sur 8.} \\
\texttt{Veuillez réessayer !} \\
            }
        }
    }
    & \\ \\

    \theutterance \stepcounter{utterance}  
    & & & \multicolumn{4}{p{0.6\linewidth}}{
        \cellcolor[rgb]{0.9,0.9,0.9}{
            \makecell[{{p{\linewidth}}}]{
                \texttt{\tiny{[GM$\langle$P2]}}
                \texttt{DIRE: Désolé, (6,5) contient '│' sur mon plateau, ce n'est pas vide. Parmi les coordonnées que vous aviez listées comme vides, pouvez{-}vous m'en proposer une autre pour C, par exemple (7,5) ou (6,6) ? Je vérifierai immédiatement si elle est libre chez moi.} \\
            }
        }
    }
     \\ \\

    \theutterance \stepcounter{utterance}  
    & & & \multicolumn{2}{p{0.3\linewidth}}{
        \cellcolor[rgb]{0.9,0.9,0.9}{
            \makecell[{{p{\linewidth}}}]{
                \texttt{\tiny{[GM$|$GM]}}
                \texttt{DIRE: Désolé, (6,5) contient '│' sur mon plateau, ce n'est pas vide. Parmi les coordonnées que vous aviez listées comme vides, pouvez{-}vous m'en proposer une autre pour C, par exemple (7,5) ou (6,6) ? Je vérifierai immédiatement si elle est libre chez moi.} \\
            }
        }
    }
    & & \\ \\

    \theutterance \stepcounter{utterance}  
    & & \multicolumn{4}{p{0.6\linewidth}}{
        \cellcolor[rgb]{0.9,0.9,0.9}{
            \makecell[{{p{\linewidth}}}]{
                \texttt{\tiny{[P1$\langle$GM]}}
                \texttt{Votre message a été transmis à l'autre joueur.} \\
\\ 
\\ 
\texttt{Vous jouez actuellement le tour 8 sur un maximum de 12.} \\
\texttt{Vous avez collectivement accumulé 7 pénalités sur 8.} \\
\texttt{L'autre joueur a envoyé ce message :} \\
\\ 
\\ 
\texttt{"Désolé, (6,5) contient '│' sur mon plateau, ce n'est pas vide. Parmi les coordonnées que vous aviez listées comme vides, pouvez{-}vous m'en proposer une autre pour C, par exemple (7,5) ou (6,6) ? Je vérifierai immédiatement si elle est libre chez moi."} \\
\\ 
\\ 
\texttt{Quelle est votre prochaine commande ?} \\
            }
        }
    }
    & \\ \\

    \theutterance \stepcounter{utterance}  
    & \multicolumn{4}{p{0.6\linewidth}}{
        \cellcolor[rgb]{0.9,0.9,0.9}{
            \makecell[{{p{\linewidth}}}]{
                \texttt{\tiny{[P1$\rangle$GM]}}
                \texttt{DIRE: D'accord, essayons (6,6) pour C. Chez moi, (6,6) est vide. Pouvez{-}vous confirmer que (6,6) est libre chez vous et y déplacer votre C ? En attendant, je vais déplacer mon P vers (3,4) ce tour.} \\
            }
        }
    }
    & & \\ \\

    \theutterance \stepcounter{utterance}  
    & & & \multicolumn{2}{p{0.3\linewidth}}{
        \cellcolor[rgb]{0.9,0.9,0.9}{
            \makecell[{{p{\linewidth}}}]{
                \texttt{\tiny{[GM$|$GM]}}
                \texttt{DIRE: D'accord, essayons (6,6) pour C. Chez moi, (6,6) est vide. Pouvez{-}vous confirmer que (6,6) est libre chez vous et y déplacer votre C ? En attendant, je vais déplacer mon P vers (3,4) ce tour.} \\
            }
        }
    }
    & & \\ \\

    \theutterance \stepcounter{utterance}  
    & & \multicolumn{4}{p{0.6\linewidth}}{
        \cellcolor[rgb]{0.9,0.9,0.9}{
            \makecell[{{p{\linewidth}}}]{
                \texttt{\tiny{[GM$\rangle$P2]}}
                \texttt{Votre message a été transmis à l'autre joueur.} \\
\\ 
\\ 
\texttt{Vous jouez actuellement le tour 9 sur un maximum de 12.} \\
\texttt{Vous avez collectivement accumulé 7 pénalités sur 8.} \\
\texttt{L'autre joueur a envoyé ce message :} \\
\\ 
\\ 
\texttt{"D'accord, essayons (6,6) pour C. Chez moi, (6,6) est vide. Pouvez{-}vous confirmer que (6,6) est libre chez vous et y déplacer votre C ? En attendant, je vais déplacer mon P vers (3,4) ce tour."} \\
\\ 
\\ 
\texttt{Quelle est votre prochaine commande ?} \\
            }
        }
    }
    & \\ \\

    \theutterance \stepcounter{utterance}  
    & & & \multicolumn{4}{p{0.6\linewidth}}{
        \cellcolor[rgb]{0.9,0.9,0.9}{
            \makecell[{{p{\linewidth}}}]{
                \texttt{\tiny{[GM$\langle$P2]}}
                \texttt{DÉPLACER: C, (6, 6)} \\
            }
        }
    }
     \\ \\

    \theutterance \stepcounter{utterance}  
    & & & \multicolumn{2}{p{0.3\linewidth}}{
        \cellcolor[rgb]{0.9,0.9,0.9}{
            \makecell[{{p{\linewidth}}}]{
                \texttt{\tiny{[GM$|$GM]}}
                \texttt{DÉPLACER: C, (6, 6)} \\
            }
        }
    }
    & & \\ \\

    \theutterance \stepcounter{utterance}  
    & & & \multicolumn{2}{p{0.3\linewidth}}{
        \cellcolor[rgb]{0.9,0.9,0.9}{
            \makecell[{{p{\linewidth}}}]{
                \texttt{\tiny{[GM$|$GM]}}
                \texttt{Pénalité : (6,6) n'est pas vide, mais contient '┴'.} \\
\\ 
\texttt{Vous avez collectivement accumulé 8 pénalités sur 8.} \\
\texttt{Veuillez réessayer !} \\
            }
        }
    }
    & & \\ \\

    \theutterance \stepcounter{utterance}  
    & & \multicolumn{4}{p{0.6\linewidth}}{
        \cellcolor[rgb]{0.9,0.9,0.9}{
            \makecell[{{p{\linewidth}}}]{
                \texttt{\tiny{[GM$\rangle$P2]}}
                \texttt{Pénalité : (6,6) n'est pas vide, mais contient '┴'.} \\
\\ 
\texttt{Vous avez collectivement accumulé 8 pénalités sur 8.} \\
\texttt{Veuillez réessayer !} \\
            }
        }
    }
    & \\ \\

    \theutterance \stepcounter{utterance}  
    & & & \multicolumn{4}{p{0.6\linewidth}}{
        \cellcolor[rgb]{0.9,0.9,0.9}{
            \makecell[{{p{\linewidth}}}]{
                \texttt{\tiny{[GM$\langle$P2]}}
                \texttt{DIRE: (6,6) contient '┴' chez moi aussi. Nous sommes à 8 pénalités, il faut faire très attention. Mes cases vides confirmées sont: (5,4) {[}L déjà placé{]}, (3,4) {[}P déjà placé{]}, (7,4), (7,5), (7,6), (5,3) et (1,1). Pouvez{-}vous vérifier si (7,5) ou (7,4) sont vides chez vous pour placer C?} \\
            }
        }
    }
     \\ \\

    \theutterance \stepcounter{utterance}  
    & & & \multicolumn{2}{p{0.3\linewidth}}{
        \cellcolor[rgb]{0.9,0.9,0.9}{
            \makecell[{{p{\linewidth}}}]{
                \texttt{\tiny{[GM$|$GM]}}
                \texttt{DIRE: (6,6) contient '┴' chez moi aussi. Nous sommes à 8 pénalités, il faut faire très attention. Mes cases vides confirmées sont: (5,4) {[}L déjà placé{]}, (3,4) {[}P déjà placé{]}, (7,4), (7,5), (7,6), (5,3) et (1,1). Pouvez{-}vous vérifier si (7,5) ou (7,4) sont vides chez vous pour placer C?} \\
            }
        }
    }
    & & \\ \\

    \theutterance \stepcounter{utterance}  
    & & \multicolumn{4}{p{0.6\linewidth}}{
        \cellcolor[rgb]{0.9,0.9,0.9}{
            \makecell[{{p{\linewidth}}}]{
                \texttt{\tiny{[P1$\langle$GM]}}
                \texttt{Votre message a été transmis à l'autre joueur.} \\
\\ 
\\ 
\texttt{Vous jouez actuellement le tour 9 sur un maximum de 12.} \\
\texttt{Vous avez collectivement accumulé 8 pénalités sur 8.} \\
\texttt{L'autre joueur a envoyé ce message :} \\
\\ 
\\ 
\texttt{"(6,6) contient '┴' chez moi aussi. Nous sommes à 8 pénalités, il faut faire très attention. Mes cases vides confirmées sont: (5,4) {[}L déjà placé{]}, (3,4) {[}P déjà placé{]}, (7,4), (7,5), (7,6), (5,3) et (1,1). Pouvez{-}vous vérifier si (7,5) ou (7,4) sont vides chez vous pour placer C?"} \\
\\ 
\\ 
\texttt{Quelle est votre prochaine commande ?} \\
            }
        }
    }
    & \\ \\

    \theutterance \stepcounter{utterance}  
    & \multicolumn{4}{p{0.6\linewidth}}{
        \cellcolor[rgb]{0.9,0.9,0.9}{
            \makecell[{{p{\linewidth}}}]{
                \texttt{\tiny{[P1$\rangle$GM]}}
                \texttt{DIRE: Sur mon plateau, (7,5) est vide. Je propose donc la configuration finale suivante : P(3,4), L(5,4), C(7,5). Je déplace mon P vers (3,4) ce tour. Pouvez{-}vous déplacer votre C vers (7,5) ce tour et confirmer ?} \\
            }
        }
    }
    & & \\ \\

    \theutterance \stepcounter{utterance}  
    & & & \multicolumn{2}{p{0.3\linewidth}}{
        \cellcolor[rgb]{0.9,0.9,0.9}{
            \makecell[{{p{\linewidth}}}]{
                \texttt{\tiny{[GM$|$GM]}}
                \texttt{DIRE: Sur mon plateau, (7,5) est vide. Je propose donc la configuration finale suivante : P(3,4), L(5,4), C(7,5). Je déplace mon P vers (3,4) ce tour. Pouvez{-}vous déplacer votre C vers (7,5) ce tour et confirmer ?} \\
            }
        }
    }
    & & \\ \\

    \theutterance \stepcounter{utterance}  
    & & \multicolumn{4}{p{0.6\linewidth}}{
        \cellcolor[rgb]{0.9,0.9,0.9}{
            \makecell[{{p{\linewidth}}}]{
                \texttt{\tiny{[GM$\rangle$P2]}}
                \texttt{Votre message a été transmis à l'autre joueur.} \\
\\ 
\\ 
\texttt{Vous jouez actuellement le tour 10 sur un maximum de 12.} \\
\texttt{Vous avez collectivement accumulé 8 pénalités sur 8.} \\
\texttt{L'autre joueur a envoyé ce message :} \\
\\ 
\\ 
\texttt{"Sur mon plateau, (7,5) est vide. Je propose donc la configuration finale suivante : P(3,4), L(5,4), C(7,5). Je déplace mon P vers (3,4) ce tour. Pouvez{-}vous déplacer votre C vers (7,5) ce tour et confirmer ?"} \\
\\ 
\\ 
\texttt{Quelle est votre prochaine commande ?} \\
            }
        }
    }
    & \\ \\

    \theutterance \stepcounter{utterance}  
    & & & \multicolumn{4}{p{0.6\linewidth}}{
        \cellcolor[rgb]{0.9,0.9,0.9}{
            \makecell[{{p{\linewidth}}}]{
                \texttt{\tiny{[GM$\langle$P2]}}
                \texttt{DÉPLACER: C, (7, 5)} \\
            }
        }
    }
     \\ \\

    \theutterance \stepcounter{utterance}  
    & & & \multicolumn{2}{p{0.3\linewidth}}{
        \cellcolor[rgb]{0.9,0.9,0.9}{
            \makecell[{{p{\linewidth}}}]{
                \texttt{\tiny{[GM$|$GM]}}
                \texttt{DÉPLACER: C, (7, 5)} \\
            }
        }
    }
    & & \\ \\

    \theutterance \stepcounter{utterance}  
    & & & \multicolumn{2}{p{0.3\linewidth}}{
        \cellcolor[rgb]{0.9,0.9,0.9}{
            \makecell[{{p{\linewidth}}}]{
                \texttt{\tiny{[GM$|$GM]}}
                \texttt{C moved to (7, 5)} \\
            }
        }
    }
    & & \\ \\

    \theutterance \stepcounter{utterance}  
    & & \multicolumn{4}{p{0.6\linewidth}}{
        \cellcolor[rgb]{0.9,0.9,0.9}{
            \makecell[{{p{\linewidth}}}]{
                \texttt{\tiny{[P1$\langle$GM]}}
                \texttt{Votre message a été transmis à l'autre joueur.} \\
\\ 
\\ 
\texttt{Vous jouez actuellement le tour 10 sur un maximum de 12.} \\
\texttt{Vous avez collectivement accumulé 8 pénalités sur 8.} \\
\texttt{L'autre joueur a déplacé un objet sur sa grille.} \\
\\ 
\\ 
\texttt{Quelle est votre prochaine commande ?} \\
            }
        }
    }
    & \\ \\

    \theutterance \stepcounter{utterance}  
    & \multicolumn{4}{p{0.6\linewidth}}{
        \cellcolor[rgb]{0.9,0.9,0.9}{
            \makecell[{{p{\linewidth}}}]{
                \texttt{\tiny{[P1$\rangle$GM]}}
                \texttt{DÉPLACER: P, (3, 4)} \\
            }
        }
    }
    & & \\ \\

    \theutterance \stepcounter{utterance}  
    & & & \multicolumn{2}{p{0.3\linewidth}}{
        \cellcolor[rgb]{0.9,0.9,0.9}{
            \makecell[{{p{\linewidth}}}]{
                \texttt{\tiny{[GM$|$GM]}}
                \texttt{DÉPLACER: P, (3, 4)} \\
            }
        }
    }
    & & \\ \\

    \theutterance \stepcounter{utterance}  
    & & & \multicolumn{2}{p{0.3\linewidth}}{
        \cellcolor[rgb]{0.9,0.9,0.9}{
            \makecell[{{p{\linewidth}}}]{
                \texttt{\tiny{[GM$|$GM]}}
                \texttt{P moved to (3, 4)} \\
            }
        }
    }
    & & \\ \\

    \theutterance \stepcounter{utterance}  
    & & \multicolumn{4}{p{0.6\linewidth}}{
        \cellcolor[rgb]{0.9,0.9,0.9}{
            \makecell[{{p{\linewidth}}}]{
                \texttt{\tiny{[GM$\rangle$P2]}}
                \texttt{Déplacement de 'C' vers (7,5) réussi. Votre grille mise à jour ressemble à ceci :} \\
\texttt{```} \\
\texttt{ 1234567} \\
\texttt{╔═╤═╤═══╗} \\
\texttt{║◌│◌│◌◌◌║ 1} \\
\texttt{║◌├─┤◌◌◌║ 2} \\
\texttt{╟─┤◌├─┬─╢ 3} \\
\texttt{║◌│P│L│◌║ 4} \\
\texttt{╟─┼─┘◌│C║ 5} \\
\texttt{║◌│◌◌┌┴─╢ 6} \\
\texttt{║◌│◌◌│◌◌║ 7} \\
\texttt{╚═╧══╧══╝} \\
\texttt{```} \\
\texttt{Vous jouez actuellement le tour 11 sur un maximum de 12.} \\
\texttt{Vous avez collectivement accumulé 8 pénalités sur 8.} \\
\texttt{L'autre joueur a déplacé un objet sur sa grille.} \\
\\ 
\\ 
\texttt{Quelle est votre prochaine commande ?} \\
            }
        }
    }
    & \\ \\

    \theutterance \stepcounter{utterance}  
    & & & \multicolumn{4}{p{0.6\linewidth}}{
        \cellcolor[rgb]{0.9,0.9,0.9}{
            \makecell[{{p{\linewidth}}}]{
                \texttt{\tiny{[GM$\langle$P2]}}
                \texttt{DIRE: Mes trois objets sont maintenant en place : P(3,4), L(5,4), C(7,5). Pouvez{-}vous confirmer que vos objets sont aux mêmes coordonnées ? Si oui, je proposerai de terminer au tour suivant.} \\
            }
        }
    }
     \\ \\

    \theutterance \stepcounter{utterance}  
    & & & \multicolumn{2}{p{0.3\linewidth}}{
        \cellcolor[rgb]{0.9,0.9,0.9}{
            \makecell[{{p{\linewidth}}}]{
                \texttt{\tiny{[GM$|$GM]}}
                \texttt{DIRE: Mes trois objets sont maintenant en place : P(3,4), L(5,4), C(7,5). Pouvez{-}vous confirmer que vos objets sont aux mêmes coordonnées ? Si oui, je proposerai de terminer au tour suivant.} \\
            }
        }
    }
    & & \\ \\

    \theutterance \stepcounter{utterance}  
    & & \multicolumn{4}{p{0.6\linewidth}}{
        \cellcolor[rgb]{0.9,0.9,0.9}{
            \makecell[{{p{\linewidth}}}]{
                \texttt{\tiny{[P1$\langle$GM]}}
                \texttt{Déplacement de 'P' vers (3,4) réussi. Votre grille mise à jour ressemble à ceci :} \\
\texttt{```} \\
\texttt{ 1234567} \\
\texttt{╔═╤═╤═══╗} \\
\texttt{║◌│◌│◌◌◌║ 1} \\
\texttt{║C├─┤◌◌◌║ 2} \\
\texttt{╟─┤◌├─┬─╢ 3} \\
\texttt{║◌│P│L│◌║ 4} \\
\texttt{╟─┼─┘◌│◌║ 5} \\
\texttt{║◌│◌◌┌┴─╢ 6} \\
\texttt{║◌│◌◌│◌◌║ 7} \\
\texttt{╚═╧══╧══╝} \\
\texttt{```} \\
\texttt{Vous jouez actuellement le tour 11 sur un maximum de 12.} \\
\texttt{Vous avez collectivement accumulé 8 pénalités sur 8.} \\
\texttt{L'autre joueur a envoyé ce message :} \\
\\ 
\\ 
\texttt{"Mes trois objets sont maintenant en place : P(3,4), L(5,4), C(7,5). Pouvez{-}vous confirmer que vos objets sont aux mêmes coordonnées ? Si oui, je proposerai de terminer au tour suivant."} \\
\\ 
\\ 
\texttt{Quelle est votre prochaine commande ?} \\
            }
        }
    }
    & \\ \\

    \theutterance \stepcounter{utterance}  
    & \multicolumn{4}{p{0.6\linewidth}}{
        \cellcolor[rgb]{0.9,0.9,0.9}{
            \makecell[{{p{\linewidth}}}]{
                \texttt{\tiny{[P1$\rangle$GM]}}
                \texttt{DÉPLACER: C, (7, 5)} \\
            }
        }
    }
    & & \\ \\

    \theutterance \stepcounter{utterance}  
    & & & \multicolumn{2}{p{0.3\linewidth}}{
        \cellcolor[rgb]{0.9,0.9,0.9}{
            \makecell[{{p{\linewidth}}}]{
                \texttt{\tiny{[GM$|$GM]}}
                \texttt{DÉPLACER: C, (7, 5)} \\
            }
        }
    }
    & & \\ \\

    \theutterance \stepcounter{utterance}  
    & & & \multicolumn{2}{p{0.3\linewidth}}{
        \cellcolor[rgb]{0.9,0.9,0.9}{
            \makecell[{{p{\linewidth}}}]{
                \texttt{\tiny{[GM$|$GM]}}
                \texttt{C moved to (7, 5)} \\
            }
        }
    }
    & & \\ \\

    \theutterance \stepcounter{utterance}  
    & & \multicolumn{4}{p{0.6\linewidth}}{
        \cellcolor[rgb]{0.9,0.9,0.9}{
            \makecell[{{p{\linewidth}}}]{
                \texttt{\tiny{[GM$\rangle$P2]}}
                \texttt{Votre message a été transmis à l'autre joueur.} \\
\\ 
\\ 
\texttt{Vous jouez actuellement le tour 12 sur un maximum de 12.} \\
\texttt{Vous avez collectivement accumulé 8 pénalités sur 8.} \\
\texttt{L'autre joueur a déplacé un objet sur sa grille.} \\
\\ 
\\ 
\texttt{Quelle est votre prochaine commande ?} \\
            }
        }
    }
    & \\ \\

    \theutterance \stepcounter{utterance}  
    & & & \multicolumn{4}{p{0.6\linewidth}}{
        \cellcolor[rgb]{0.9,0.9,0.9}{
            \makecell[{{p{\linewidth}}}]{
                \texttt{\tiny{[GM$\langle$P2]}}
                \texttt{DIRE: terminé ?} \\
            }
        }
    }
     \\ \\

    \theutterance \stepcounter{utterance}  
    & & & \multicolumn{2}{p{0.3\linewidth}}{
        \cellcolor[rgb]{0.9,0.9,0.9}{
            \makecell[{{p{\linewidth}}}]{
                \texttt{\tiny{[GM$|$GM]}}
                \texttt{DIRE: terminé ?} \\
            }
        }
    }
    & & \\ \\

    \theutterance \stepcounter{utterance}  
    & & & \multicolumn{2}{p{0.3\linewidth}}{
        \cellcolor[rgb]{0.9,0.9,0.9}{
            \makecell[{{p{\linewidth}}}]{
                \texttt{\tiny{[GM$|$GM]}}
                \texttt{Maximum number of rounds reached} \\
            }
        }
    }
    & & \\ \\

    \theutterance \stepcounter{utterance}  
    & & & \multicolumn{2}{p{0.3\linewidth}}{
        \cellcolor[rgb]{0.9,0.9,0.9}{
            \makecell[{{p{\linewidth}}}]{
                \texttt{\tiny{[GM$|$GM]}}
                \texttt{Initial states:} \\
\texttt{```} \\
\texttt{Player 1:} \\
\texttt{ 1234567} \\
\texttt{╔═╤═╤═══╗} \\
\texttt{║L│◌│◌◌◌║ 1} \\
\texttt{║C├─┤◌◌◌║ 2} \\
\texttt{╟─┤◌├─┬─╢ 3} \\
\texttt{║◌│◌│◌│◌║ 4} \\
\texttt{╟─┼─┘◌│◌║ 5} \\
\texttt{║◌│◌◌┌┴─╢ 6} \\
\texttt{║◌│◌◌│◌P║ 7} \\
\texttt{╚═╧══╧══╝} \\
\\ 
\\ 
\texttt{Player 2:} \\
\texttt{ 1234567} \\
\texttt{╔═╤═╤═══╗} \\
\texttt{║◌│◌│◌◌◌║ 1} \\
\texttt{║◌├─┤◌L◌║ 2} \\
\texttt{╟─┤P├─┬─╢ 3} \\
\texttt{║◌│◌│◌│◌║ 4} \\
\texttt{╟─┼─┘C│◌║ 5} \\
\texttt{║◌│◌◌┌┴─╢ 6} \\
\texttt{║◌│◌◌│◌◌║ 7} \\
\texttt{╚═╧══╧══╝} \\
\\ 
\texttt{```} \\
            }
        }
    }
    & & \\ \\

    \theutterance \stepcounter{utterance}  
    & & & \multicolumn{2}{p{0.3\linewidth}}{
        \cellcolor[rgb]{0.9,0.9,0.9}{
            \makecell[{{p{\linewidth}}}]{
                \texttt{\tiny{[GM$|$GM]}}
                \texttt{End states:} \\
\texttt{```} \\
\texttt{Player 1:} \\
\texttt{ 1234567} \\
\texttt{╔═╤═╤═══╗} \\
\texttt{║◌│◌│◌◌◌║ 1} \\
\texttt{║◌├─┤◌◌◌║ 2} \\
\texttt{╟─┤◌├─┬─╢ 3} \\
\texttt{║◌│P│L│◌║ 4} \\
\texttt{╟─┼─┘◌│C║ 5} \\
\texttt{║◌│◌◌┌┴─╢ 6} \\
\texttt{║◌│◌◌│◌◌║ 7} \\
\texttt{╚═╧══╧══╝} \\
\\ 
\\ 
\texttt{Player 2:} \\
\texttt{ 1234567} \\
\texttt{╔═╤═╤═══╗} \\
\texttt{║◌│◌│◌◌◌║ 1} \\
\texttt{║◌├─┤◌◌◌║ 2} \\
\texttt{╟─┤◌├─┬─╢ 3} \\
\texttt{║◌│P│L│◌║ 4} \\
\texttt{╟─┼─┘◌│C║ 5} \\
\texttt{║◌│◌◌┌┴─╢ 6} \\
\texttt{║◌│◌◌│◌◌║ 7} \\
\texttt{╚═╧══╧══╝} \\
\\ 
\texttt{```} \\
            }
        }
    }
    & & \\ \\

    \theutterance \stepcounter{utterance}  
    & & & \multicolumn{2}{p{0.3\linewidth}}{
        \cellcolor[rgb]{0.9,0.9,0.9}{
            \makecell[{{p{\linewidth}}}]{
                \texttt{\tiny{[GM$|$GM]}}
                \texttt{* success: True} \\
\texttt{* lose: False} \\
\texttt{* aborted: False} \\
\texttt{{-}{-}{-}{-}{-}{-}{-}} \\
\texttt{* Shifts: 4.00} \\
\texttt{* Max Shifts: 4.00} \\
\texttt{* Min Shifts: 2.00} \\
\texttt{* End Distance Sum: 0.00} \\
\texttt{* Init Distance Sum: 15.76} \\
\texttt{* Expected Distance Sum: 12.57} \\
\texttt{* Penalties: 8.00} \\
\texttt{* Max Penalties: 8.00} \\
\texttt{* Rounds: 12.00} \\
\texttt{* Max Rounds: 12.00} \\
\texttt{* Object Count: 3.00} \\
            }
        }
    }
    & & \\ \\

\end{supertabular}
}

\end{document}
